
\chapter*{Exposé XV\\Complements on Sub-tori of a Group Prescheme\\
Application to Smooth Groups}
\addcontentsline{toc}{chapter}{Exposé XV. Complements on Sub-tori of a Group Prescheme}
\SGARefTarget{sga3:xv:expose}

\begin{center}
\textit{by M. Raynaud}\textsuperscript{,}\footnote{This Exposé and the two following ones
(Exposés XVI and XVII) do not correspond to oral lectures delivered at
the seminar. With some additional material, they develop the substance
of concise manuscript notes written by A. Grothendieck in 1964 for the
present seminar.}
\end{center}

\SGASection{0}{Introduction}
\SGARefTarget{sga3:xv:section:0}

This Exposé supplements and partly recasts Exposés XI and XII; knowledge
of Exposés XIII and XIV is not indispensable. Continuing the effort
begun in XII, we shall work with \(S\)-group preschemes that are not
necessarily affine and not necessarily separated over \(S\).

\SGARefLink{sga3:xv:section:1}{Sections 1}, 2, 3, and 4 are devoted to the study of sub-tori of a group
prescheme. We obtain infinitesimal lifting theorems (§2) and a global
lifting theorem (§4), in which points of finite order (§1) play an
essential role.

\SGARefLink{sga3:xv:section:5}{Sections 5}, 6, and 7 are independent of the preceding sections. The
consideration of infinitesimal neighbourhoods leads to the
representability of the functor of smooth subgroups equal to their
connected normalizer (§5). In §§6 and 7 we turn more particularly to
Cartan subgroups.

Finally, in §8 we give a necessary and sufficient condition for the
functor of sub-tori of a smooth group, or the functor of maximal tori,
to be representable.
